\documentclass[a4paper]{arrowhead}

\usepackage[yyyymmdd]{datetime}
\usepackage{etoolbox}
\usepackage[utf8]{inputenc}
\usepackage{multirow}

\renewcommand{\dateseparator}{-}

%% Special references
\newcommand{\fref}[1]{{\textcolor{ArrowheadBlue}{\hyperref[sec:functions:#1]{#1}}}}
\newcommand{\mref}[1]{{\textcolor{ArrowheadPurple}{\hyperref[sec:model:#1]{#1}}}}
\newcommand{\pdef}[1]{{\textcolor{ArrowheadGrey}{#1\label{sec:model:primitives:#1}\label{sec:model:primitives:#1s}\label{sec:model:primitives:#1es}}}}
\newcommand{\pref}[1]{{\textcolor{ArrowheadGrey}{\hyperref[sec:model:primitives:#1]{#1}}}}

\newrobustcmd\fsubsection[3]{
  \addtocounter{subsection}{1}
  \addcontentsline{toc}{subsection}{\protect\numberline{\thesubsection}function \textcolor{ArrowheadBlue}{#1}}
  \renewcommand*{\do}[1]{\rref{##1},\ }
  \subsection*{
    \thesubsection\quad
    operation
    \textcolor{ArrowheadBlue}{#1}
    (\notblank{#2}{\mref{#2}}{})
    \notblank{#3}{: \mref{#3}}{}
  }
  \label{sec:functions:#1}
}
\newrobustcmd\msubsection[2]{
  \addtocounter{subsection}{1}
  \addcontentsline{toc}{subsection}{\protect\numberline{\thesubsection}#1 \textcolor{ArrowheadPurple}{#2}}
  \subsection*{\thesubsection\quad#1 \textcolor{ArrowheadPurple}{#2}}
  \label{sec:model:#2} \label{sec:model:#2s} \label{sec:model:#2es}
}
%%

\begin{document}

%% Arrowhead Document Properties
\ArrowheadTitle{Eclispe Arrohead Naming Convention} % XXX = SystemName e.g. Service Registry HTTP/TLS/JSON}
\ArrowheadType{}
\ArrowheadTypeShort{SoSD}
\ArrowheadVersion{5.0.0} % Arrowhead version X.Y.Z, e..g. 4.4.1
\ArrowheadDate{\today}
\ArrowheadAuthor{Jerker Delsing} % Corresponding author e.g. Jerker Delsing
\ArrowheadStatus{Prototype} % e..g. RELEASE, RELEASE CONDIDATE, PROTOTYPE
\ArrowheadContact{jerker.delsing@ltu.se} % Email of corresponding author
\ArrowheadFooter{\href{www.arrowhead.eu}{www.arrowhead.eu}}
\ArrowheadSetup
%%

%% Front Page
\begin{center}
  \vspace*{1cm}
  \huge{\arrowtitle}

  \vspace*{0.2cm}
  \LARGE{\arrowtype}
  \vspace*{1cm}

  %\Large{Service ID: \textit{"\arrowid"}}
  \vspace*{\fill}

  % Front Page Image
  %\includegraphics{figures/TODO}

  \vspace*{1cm}
  \vspace*{\fill}

  % Front Page Abstract
  \begin{abstract}
    Proposla for naming convention of microsystems, microservices and
    associated attributes and metadata. This is intended as an
    appendix to the Eclipse Arrowhead GSoSD document.  
  \end{abstract}

  \vspace*{1cm}

%   \scriptsize
%   \begin{tabularx}{\textwidth}{l X}
%     \raisebox{-0.5\height}{\includegraphics[width=2cm]{figures/artemis_logo}} & {ARTEMIS Innovation Pilot Project: Arrowhead\newline
%     THEME [SP1-JTI-ARTEMIS-2012-AIPP4 SP1-JTI-ARTEMIS-2012-AIPP6]\newline
%     [Production and Energy System Automation Intelligent-Built environment and urban infrastructure for sustainable and friendly cities]}
%   \end{tabularx}
%   \vspace*{-0.2cm}
 \end{center}

\newpage
%%

%% Table of Contents
\tableofcontents
\newpage
%%

\section{Overview}
\label{sec:overview}
  Proposla for naming convention of microsystems, microservices and
    associated attributes and metadata. This is intended as an
    appendix to the Eclipse Arrowhead GSoSD document.
    
The rest of this document is organized as follows.
In Section \ref{sec:prior_art}, we reference major prior art on
microsystem and microservice naming within the Eclispe Arrowhead project. 

In Section \ref{sec:principles}, we detail the underlaying thinking
and principles for the
naming convention.

In Section \ref{sec:examples}, we provide a set of example.


\newpage

\subsection{Significant Prior Art}
\label{sec:prior_art}

A previous proposal by Paniagua et.al \cite{Paniagua_2019} has not
  gained attention. Thus we here proposa a significantly simpler
  naming convention approach for Eclipse Arrowhead mincrosystems,
  microservice and associated metadata and attributes.

\subsection{Foundational naming principles}
\label{sec:principles}
The ambition with this naming conventions is to provide names being:
\begin{itemize}
\item Avoid namings which generates considerable extra work for most programing
  languages
\item Names shall reflect the intended functionality and usage in an SOA
  architecture
\item Microsystem and Microservice name shall not be identical
\item Microsystems name shall start with a Capital letter
\item Microservices name shall starrts with a lowercase letter

\item Metadata and attribute naming should be descriptive enabling
  easy human understandning of context
  
\item Naming of instances of microsystems microservices, metadata and
  attrtibutes shall follow  the naming convention of the
  applied standard
  
\item The choice of indsutry standards to be applied should preferable
  be possible to connected to the Industrial Data Ontology (IDO), ISO 23726-3
\end{itemize}


\subsection{Naming example}
\label{sec:examples}

PLease find below a set of examples for the most important naming
siutaitons in the Eclipse Arrowhead architecture
\begin{itemize}
\item Microsystems name always start with a Capital letter: examples
  \begin{itemize}
  \item ServiceRegistry
  \item DynamicServiceOrchestration
  \item SimpleServiceOrchestration
  \item FlexibleServiceOrchestration

  \item ComputeOrchestrationSystem
  \item DeploymentOrchestrationSystem
  \item ConsumerAuhtorizationSystem
  \item Authentication
  \end{itemize}
  
\item Microservices name always starrts with a lowercase letter: examples 
  \begin{itemize}
  \item serviceDiscovery
  \item serviceOrchestration
  \item computeOrchestration
  \item simpleOrchestrationStoreManagement
  \item flexibleOrchestrationStoreManagement
  \item deploymentOrchestration
  \item consumerAuhtorization
  \end{itemize}

\item Metadata and attribute naming which in combination to the
  related Microsystem or Micsroservice shall be meaningfull
  \begin{itemize}
  \item Metadata/attribute: Timestamp (of what should be possible to infere from
    the naming of the microsystem or microservice instance to which the
    metadata/attribute is connected
  \item Metadata/attribute: SWversion (version reference of the
    deployed software)
  \item Metadata/attribute: SWcompiler (which comiler was used)
  \item Metadata/attribute: CompilerSwitches (used compiler switches
    and value) 
  \end{itemize}

\item Instance naming shall follow the naming convention of the
    applied standard e.g. ISO 15296, ISO 10303, S5000. 

\end{itemize} 


\bibliographystyle{IEEEtran}
\bibliography{bibliography}

\newpage

\section{Revision History}
%\subsection{Amendments}


\noindent\begin{tabularx}{\textwidth}{| p{1cm} | p{3cm} | p{2cm} | X | p{4cm} |} \hline
\rowcolor{gray!33} No. & Date & Version & Subject of Amendments & Author \\ \hline

1 & 2025-04-02 & \arrowversion & & Jerker Delsing \\ \hline
2 & 2025-04-09 & \arrowversion & & Jerker Delsing \\ \hline
3 & & & & \\ \hline

\end{tabularx}

\subsection{Quality Assurance}

\noindent\begin{tabularx}{\textwidth}{| p{1cm} | p{3cm} | p{2cm} | X |} \hline
\rowcolor{gray!33} No. & Date & Version & Approved by \\ \hline

1 & 2025-04-02 & \arrowversion  &  Pal Varga\\ \hline

\end{tabularx}

\end{document}